\uuid{ACHS}
\exo7id{4035}
\titre{exo7 4035}
\auteur{quercia}
\organisation{exo7}
\datecreate{2010-03-11}
\isIndication{false}
\isCorrection{false}
\chapitre{Développement limité}
\sousChapitre{Autre}
\cpge{mpsi-mp2i.convexite.b, mpsi-mp2i.analyse-asymptotique.b, pcsi.analyse-asymptotique.b, ptsi.analyse-asymptotique.b, mpsi-mp2i.derivabilite.a, pcsi.derivabilite.a, ptsi.derivabilite.a, bcpst1.developpements-limites.a, bcpst1.derivation.a}
\module{Analyse}
\niveau{L1}
\difficulte{}

\contenu{
\question{
Soit $f : \R \to \R$ convexe dérivable telle que
$f(a+h) = f(a) + hf'(a) + o (h^2)$.

Démontrer que $f$ est deux fois dérivable en $a$ et $f''(a) = 0$
(comparer $f'(a+h)$ aux taux d'accroissement de $f$ entre $a$ et $a+h$,
et entre $a+h$ et $a+2h$).

\'Etudier le cas où $f(a+h) = f(a) + hf'(a) + \frac {Lh^2}2 + o (h^2)$.
}
}
